\chapter{CONTRIBUIÇÕES TEÓRICAS E PRÁTICAS ESPERADAS}

Espera-se que esta pesquisa contribua para o debate sobre a educação geográfica, ao articular os conceitos de espaço e cidadania. Do ponto de vista teórico, o trabalho busca fortalecer o diálogo entre a geografia crítica de \cite{milton-ne} e as práticas pedagógicas contemporâneas defendidas por \cite{tec-ensi-pqn}.

Na prática, o estudo pretende oferecer subsídios para professores da rede básica, sugerindo abordagens que valorizem a leitura crítica do espaço e incentivem a participação ativa dos estudantes na transformação de seu território.